% Options for packages loaded elsewhere
\PassOptionsToPackage{unicode}{hyperref}
\PassOptionsToPackage{hyphens}{url}
\documentclass[
]{article}
\usepackage{xcolor}
\usepackage{amsmath,amssymb}
\setcounter{secnumdepth}{5}
\usepackage{iftex}
\ifPDFTeX
  \usepackage[T1]{fontenc}
  \usepackage[utf8]{inputenc}
  \usepackage{textcomp} % provide euro and other symbols
\else % if luatex or xetex
  \usepackage{unicode-math} % this also loads fontspec
  \defaultfontfeatures{Scale=MatchLowercase}
  \defaultfontfeatures[\rmfamily]{Ligatures=TeX,Scale=1}
\fi
\usepackage{lmodern}
\ifPDFTeX\else
  % xetex/luatex font selection
\fi
% Use upquote if available, for straight quotes in verbatim environments
\IfFileExists{upquote.sty}{\usepackage{upquote}}{}
\IfFileExists{microtype.sty}{% use microtype if available
  \usepackage[]{microtype}
  \UseMicrotypeSet[protrusion]{basicmath} % disable protrusion for tt fonts
}{}
\makeatletter
\@ifundefined{KOMAClassName}{% if non-KOMA class
  \IfFileExists{parskip.sty}{%
    \usepackage{parskip}
  }{% else
    \setlength{\parindent}{0pt}
    \setlength{\parskip}{6pt plus 2pt minus 1pt}}
}{% if KOMA class
  \KOMAoptions{parskip=half}}
\makeatother
\usepackage{color}
\usepackage{fancyvrb}
\newcommand{\VerbBar}{|}
\newcommand{\VERB}{\Verb[commandchars=\\\{\}]}
\DefineVerbatimEnvironment{Highlighting}{Verbatim}{commandchars=\\\{\}}
% Add ',fontsize=\small' for more characters per line
\usepackage{framed}
\definecolor{shadecolor}{RGB}{248,248,248}
\newenvironment{Shaded}{\begin{snugshade}}{\end{snugshade}}
\newcommand{\AlertTok}[1]{\textcolor[rgb]{0.94,0.16,0.16}{#1}}
\newcommand{\AnnotationTok}[1]{\textcolor[rgb]{0.56,0.35,0.01}{\textbf{\textit{#1}}}}
\newcommand{\AttributeTok}[1]{\textcolor[rgb]{0.13,0.29,0.53}{#1}}
\newcommand{\BaseNTok}[1]{\textcolor[rgb]{0.00,0.00,0.81}{#1}}
\newcommand{\BuiltInTok}[1]{#1}
\newcommand{\CharTok}[1]{\textcolor[rgb]{0.31,0.60,0.02}{#1}}
\newcommand{\CommentTok}[1]{\textcolor[rgb]{0.56,0.35,0.01}{\textit{#1}}}
\newcommand{\CommentVarTok}[1]{\textcolor[rgb]{0.56,0.35,0.01}{\textbf{\textit{#1}}}}
\newcommand{\ConstantTok}[1]{\textcolor[rgb]{0.56,0.35,0.01}{#1}}
\newcommand{\ControlFlowTok}[1]{\textcolor[rgb]{0.13,0.29,0.53}{\textbf{#1}}}
\newcommand{\DataTypeTok}[1]{\textcolor[rgb]{0.13,0.29,0.53}{#1}}
\newcommand{\DecValTok}[1]{\textcolor[rgb]{0.00,0.00,0.81}{#1}}
\newcommand{\DocumentationTok}[1]{\textcolor[rgb]{0.56,0.35,0.01}{\textbf{\textit{#1}}}}
\newcommand{\ErrorTok}[1]{\textcolor[rgb]{0.64,0.00,0.00}{\textbf{#1}}}
\newcommand{\ExtensionTok}[1]{#1}
\newcommand{\FloatTok}[1]{\textcolor[rgb]{0.00,0.00,0.81}{#1}}
\newcommand{\FunctionTok}[1]{\textcolor[rgb]{0.13,0.29,0.53}{\textbf{#1}}}
\newcommand{\ImportTok}[1]{#1}
\newcommand{\InformationTok}[1]{\textcolor[rgb]{0.56,0.35,0.01}{\textbf{\textit{#1}}}}
\newcommand{\KeywordTok}[1]{\textcolor[rgb]{0.13,0.29,0.53}{\textbf{#1}}}
\newcommand{\NormalTok}[1]{#1}
\newcommand{\OperatorTok}[1]{\textcolor[rgb]{0.81,0.36,0.00}{\textbf{#1}}}
\newcommand{\OtherTok}[1]{\textcolor[rgb]{0.56,0.35,0.01}{#1}}
\newcommand{\PreprocessorTok}[1]{\textcolor[rgb]{0.56,0.35,0.01}{\textit{#1}}}
\newcommand{\RegionMarkerTok}[1]{#1}
\newcommand{\SpecialCharTok}[1]{\textcolor[rgb]{0.81,0.36,0.00}{\textbf{#1}}}
\newcommand{\SpecialStringTok}[1]{\textcolor[rgb]{0.31,0.60,0.02}{#1}}
\newcommand{\StringTok}[1]{\textcolor[rgb]{0.31,0.60,0.02}{#1}}
\newcommand{\VariableTok}[1]{\textcolor[rgb]{0.00,0.00,0.00}{#1}}
\newcommand{\VerbatimStringTok}[1]{\textcolor[rgb]{0.31,0.60,0.02}{#1}}
\newcommand{\WarningTok}[1]{\textcolor[rgb]{0.56,0.35,0.01}{\textbf{\textit{#1}}}}
\usepackage{longtable,booktabs,array}
\newcounter{none} % for unnumbered tables
\usepackage{calc} % for calculating minipage widths
% Correct order of tables after \paragraph or \subparagraph
\usepackage{etoolbox}
\makeatletter
\patchcmd\longtable{\par}{\if@noskipsec\mbox{}\fi\par}{}{}
\makeatother
% Allow footnotes in longtable head/foot
\IfFileExists{footnotehyper.sty}{\usepackage{footnotehyper}}{\usepackage{footnote}}
\makesavenoteenv{longtable}
\setlength{\emergencystretch}{3em} % prevent overfull lines
\providecommand{\tightlist}{%
  \setlength{\itemsep}{0pt}\setlength{\parskip}{0pt}}
\usepackage{bookmark}
\IfFileExists{xurl.sty}{\usepackage{xurl}}{} % add URL line breaks if available
\urlstyle{same}
\hypersetup{
  hidelinks,
  pdfcreator={LaTeX via pandoc}}

\author{}
\date{}

\begin{document}

{
\setcounter{tocdepth}{3}
\tableofcontents
}
\section{1. Overview}\label{overview}

Workwarrior is a profile-based shell productivity system that wraps five
open-source tools --- TaskWarrior, TimeWarrior, JRNL, Hledger, and
Bugwarrior --- into a single unified command surface. The system
provides:

\begin{itemize}
\tightlist
\item
  \textbf{Profile isolation} --- each profile is an independent
  directory containing all tool data for one work context
\item
  \textbf{Unified CLI} --- a single \texttt{ww} dispatcher routes all
  commands to service scripts
\item
  \textbf{Natural language input} --- 627 compiled heuristic regex rules
  translate plain English to tool commands before optionally falling
  back to a local or remote LLM
\item
  \textbf{Browser UI} --- a locally-served web interface (Python 3
  stdlib only) with 15+ panels, real-time SSE updates, and a unified
  command input
\item
  \textbf{GitHub sync} --- two complementary engines for pulling issues
  from 20+ services and bidirectional sync between individual tasks and
  GitHub issues
\item
  \textbf{Extensible service architecture} --- new services are
  executable scripts in
  \texttt{services/\textless{}category\textgreater{}/}, discovered at
  runtime with no registration step
\end{itemize}

\textbf{Entry point:} \texttt{bin/ww} (main CLI dispatcher)\\
\textbf{Shell bootstrap:} \texttt{bin/ww-init.sh} (sourced by
\texttt{.bashrc}/\texttt{.zshrc} at shell start)\\
\textbf{Canonical install path:} \texttt{\textasciitilde{}/ww}

\begin{center}\rule{0.5\linewidth}{0.5pt}\end{center}

\section{2. Architecture}\label{architecture}

\begin{verbatim}
User types: p-work
  → lib/shell-integration.sh sets TASKRC, TASKDATA, TIMEWARRIORDB, etc.
  → all tools now operate on the work profile's data

User types: task add "Ship API" due:friday
  → TaskWarrior writes to profiles/work/.task/
  → on-modify hook starts TimeWarrior tracking

User types: ww browser
  → Python3 HTTP server starts on localhost:7777
  → serves static UI + REST API + SSE events
  → reads profile data via same env vars

User types in browser CMD: "add task review code and start tracking"
  → HeuristicEngine matches against 627 rules
  → splits compound command on "and"
  → executes: task add review code + timew start review code
\end{verbatim}

\subsection{Data Flow Summary}\label{data-flow-summary}

\begin{verbatim}
Shell init
  ww-init.sh sourced
    → shell-integration.sh defines p-<name> aliases and shell functions
    → ww alias registered

Profile activation (p-work)
  → use_task_profile("work") called
  → 5 env vars exported
  → all subsequent tool calls use profile data

ww <command> [args]
  → bin/ww validates profile active (or allow-listed command)
  → discovers service script in services/<command>/
  → profile-level service checked first (overrides global)
  → executes service with remaining args

ww browser
  → services/browser/server.py starts ThreadingHTTPServer on :7777
  → SSE channel open for real-time profile events
  → REST API handles panel data and CMD execution
  → CMD input → HeuristicEngine → 627 rules → AI fallback → ww subcommand execution
\end{verbatim}

\begin{center}\rule{0.5\linewidth}{0.5pt}\end{center}

\section{3. Directory Structure}\label{directory-structure}

\begin{verbatim}
bin/
  ww                          Main CLI dispatcher (all ww commands route here)
  ww-init.sh                  Shell bootstrap (sourced at shell start)
  profile                     Profile activation helper
  custom                      Custom service launcher
  export                      Export launcher
  x                           Delete launcher

lib/                          Core bash libraries (sourced, not executed directly)
  core-utils.sh               WW_BASE resolution, profile validation, last-profile state
  profile-manager.sh          Profile lifecycle: create, delete, backup, import, restore
  shell-integration.sh        Shell function injection, alias management, rc file writes
  logging.sh                  log_info/log_success/log_warning/log_error/log_step
  github-api.sh               GitHub REST API via gh CLI
  github-sync-state.sh        Sync state persistence (JSON per task)
  sync-pull.sh                GitHub → TaskWarrior pull
  sync-push.sh                TaskWarrior → GitHub push
  sync-bidirectional.sh       Orchestrates pull+push with conflict window
  sync-detector.sh            Change detection (task side + GitHub side)
  field-mapper.sh             Field mapping between TaskWarrior and GitHub formats
  conflict-resolver.sh        Last-write-wins conflict resolution with configurable window
  annotation-sync.sh          TaskWarrior annotations ↔ GitHub comments sync
  sync-permissions.sh         Per-UDA sync permission tokens
  config-loader.sh            Profile config loading (jrnl.yaml, ledgers.yaml, bugwarrior)
  config-utils.sh             YAML parsing utilities
  taskwarrior-api.sh          tw_get_task, tw_update_task, tw_get_field wrappers
  bugwarrior-integration.sh   Bugwarrior config and pull helpers
  dependency-installer.sh     Platform-aware tool installer with version cards
  installer-utils.sh          Install path resolution, shell RC block management
  export-utils.sh             Profile data export (JSON, CSV, markdown)
  delete-utils.sh             Profile deletion with backup safety
  profile-stats.sh            Profile statistics and reporting
  shortcode-registry.sh       Shortcut/alias registry read/write
  error-handler.sh            GitHub error classification

services/                     Service scripts (executable, discovered at runtime)
  profile/
    create-ww-profile.sh      Full profile creation (dirs, taskrc, hooks, aliases)
    manage-profiles.sh        list/info/delete/backup/import/restore
    profile-tool.sh           Profile info display
    urgency.sh                ww profile urgency — urgency coefficient management
    subservices/
      profile-uda.sh          ww profile uda — full UDA management surface
      profile-density.sh      ww profile density — TWDensity integration
      uda-manager.sh          Interactive UDA CRUD
  custom/
    configure-journals.sh     ww custom journals
    configure-ledgers.sh      ww custom ledgers
    configure-tasks.sh        ww custom tasks
    configure-times.sh        ww custom times
    configure-issues.sh       ww custom issues / ww issues custom
    github-sync.sh            ww issues push/pull/sync/enable/disable/status
  questions/
    q.sh                      Main questions dispatcher
    handlers/                 Per-service question handlers
    templates/                YAML question templates
  extensions/
    extensions.sh             ww extensions taskwarrior list/search/info
    taskwarrior.py            Extension data fetcher
  find/
    find.sh                   ww find dispatcher
    find.py                   Python search implementation
  groups/
    groups.sh                 ww group list/show/create/add/remove/delete
  models/
    models.sh                 ww model list/providers/add/set-default
  export/
    export.sh                 ww export (JSON, CSV, markdown)
  browser/
    server.py                 Python3 HTTP server — 15+ panels, SSE, CMD, heuristic engine
    static/app.js             Browser UI — sidebar, task editor, all panels
    static/index.html         HTML shell — dark terminal aesthetic
    static/style.css          Styles
  cmd/
    cmd.log                   JSONL log of all CMD submissions
  ctrl/
    ctrl.sh                   ww ctrl status/ai-mode/prompt-ww/prompt-ai
  x-delete/
    x.sh                      ww x — profile/data deletion with safety backups
  servers/                    Server management
  scripts/                    Utility scripts
  find/                       Cross-profile search

weapons/
  gun/                        taskgun passthrough — bulk task series generator
  sword/                      Task splitting into sequential subtasks with dependencies

scripts/
  compile-heuristics.py       Heuristic compilation
  scan-taskwarrior-extensions.py  Extension registry scanner

functions/                    Shell helper functions (sourced at init)
  journals/                   Journal helpers
  ledgers/                    Ledger helpers
  tasks/                      TaskWarrior helpers and extensions
  times/                      TimeWarrior helpers
  issues/                     Issue tracking helpers

tools/
  list/list.py                List management tool (ww list)

config/
  ai.yaml                     AI mode, access points, preferred provider
  cmd-heuristics.yaml         Compiled NL→command rules (627 rules)
  cmd-heuristics-corpus.yaml  Synthetic corpus for heuristic validation
  models.yaml                 LLM provider/model registry
  ctrl.yaml                   CTRL service settings
  groups.yaml                 Profile group definitions
  shortcuts.yaml              Shortcut definitions

profiles/                     User profiles (created at runtime, gitignored)
  <name>/
    .taskrc                   TaskWarrior config (UDAs, reports, urgency)
    .task/                    TaskWarrior database
    .timewarrior/             TimeWarrior database
    journals/                 Journal files
    ledgers/                  Hledger ledger files
    jrnl.yaml                 Journal name → file mapping
    ledgers.yaml              Ledger name → file mapping
    .config/                  Service configs (bugwarrior)

system/                       Control plane (not shipped to users)
  CLAUDE.md                   Primary agent context document
  TASKS.md                    Canonical task board
  fragility-register.md       File fragility classifications
  gates/                      Gate contract templates
  roles/                      Agent role definitions
  plans/                      Planning documents
  specs/                      Specification documents
  audits/                     Audit outputs
  workflows/                  Process workflows
\end{verbatim}

\begin{center}\rule{0.5\linewidth}{0.5pt}\end{center}

\section{4. Environment Variables}\label{environment-variables}

Set automatically by \texttt{use\_task\_profile()} when a profile is
activated via \texttt{p-\textless{}name\textgreater{}}:

{\def\LTcaptype{none} % do not increment counter
\begin{longtable}[]{@{}
  >{\raggedright\arraybackslash}p{(\linewidth - 4\tabcolsep) * \real{0.3846}}
  >{\raggedright\arraybackslash}p{(\linewidth - 4\tabcolsep) * \real{0.2692}}
  >{\raggedright\arraybackslash}p{(\linewidth - 4\tabcolsep) * \real{0.3462}}@{}}
\toprule\noalign{}
\begin{minipage}[b]{\linewidth}\raggedright
Variable
\end{minipage} & \begin{minipage}[b]{\linewidth}\raggedright
Value
\end{minipage} & \begin{minipage}[b]{\linewidth}\raggedright
Used By
\end{minipage} \\
\midrule\noalign{}
\endhead
\bottomrule\noalign{}
\endlastfoot
\texttt{WARRIOR\_PROFILE} & Active profile name (e.g.~\texttt{work}) &
All ww scripts \\
\texttt{WORKWARRIOR\_BASE} & Profile base directory
(\texttt{\textasciitilde{}/ww/profiles/work}) & All ww scripts \\
\texttt{TASKRC} & Path to profile \texttt{.taskrc} & TaskWarrior \\
\texttt{TASKDATA} & Path to profile \texttt{.task/} & TaskWarrior \\
\texttt{TIMEWARRIORDB} & Path to profile \texttt{.timewarrior/} &
TimeWarrior \\
\end{longtable}
}

\textbf{Rule:} Nothing in the system hardcodes paths. All path
resolution goes through these variables. A lib function or service that
constructs a path by concatenating
\texttt{\textasciitilde{}/ww/profiles/\textless{}name\textgreater{}/...}
is a bug --- it must use \texttt{\$WORKWARRIOR\_BASE}.

\begin{center}\rule{0.5\linewidth}{0.5pt}\end{center}

\section{5. Profile Model}\label{profile-model}

\subsection{Structure}\label{structure}

\begin{verbatim}
profiles/<name>/
  .taskrc              TaskWarrior config (UDAs, reports, urgency coefficients, hooks)
  .task/               TaskWarrior database (taskchampion.sqlite3 + supporting files)
  .timewarrior/        TimeWarrior database (data/ + .timewarrior config)
  journals/            JRNL journal files
  ledgers/             Hledger ledger files
  jrnl.yaml            Journal name → file path mapping
  ledgers.yaml         Ledger name → file path mapping
  .config/             Per-profile service configs
    bugwarrior/        Bugwarrior config (bugwarriorrc or bugwarrior.toml)
    taskcheck/         Taskcheck config
  services/            Per-profile service overrides (shadow global services)
  ai.yaml              Per-profile AI mode override (optional)
\end{verbatim}

\subsection{Lifecycle}\label{lifecycle}

\begin{Shaded}
\begin{Highlighting}[]
\ExtensionTok{ww}\NormalTok{ profile create }\OperatorTok{\textless{}}\NormalTok{name}\OperatorTok{\textgreater{}}       \CommentTok{\# Creates directory structure, writes .taskrc, injects hooks, creates p{-}\textless{}name\textgreater{} alias}
\ExtensionTok{ww}\NormalTok{ profile list                }\CommentTok{\# Lists all profiles with last{-}active timestamps}
\ExtensionTok{ww}\NormalTok{ profile info }\OperatorTok{\textless{}}\NormalTok{name}\OperatorTok{\textgreater{}}         \CommentTok{\# Shows profile details, resource inventory, sync state}
\ExtensionTok{ww}\NormalTok{ profile delete }\OperatorTok{\textless{}}\NormalTok{name}\OperatorTok{\textgreater{}}       \CommentTok{\# Marks deleted, creates safety backup}
\ExtensionTok{ww}\NormalTok{ profile backup }\OperatorTok{\textless{}}\NormalTok{name}\OperatorTok{\textgreater{}}       \CommentTok{\# Archives to profiles/\textless{}name\textgreater{}{-}backup{-}\textless{}timestamp\textgreater{}.tar.gz}
\ExtensionTok{ww}\NormalTok{ profile import }\OperatorTok{\textless{}}\NormalTok{archive}\OperatorTok{\textgreater{}}    \CommentTok{\# Creates new profile from archive}
\ExtensionTok{ww}\NormalTok{ profile restore }\OperatorTok{\textless{}}\NormalTok{archive}\OperatorTok{\textgreater{}}   \CommentTok{\# Replaces existing profile from archive}
\ExtensionTok{ww}\NormalTok{ remove }\OperatorTok{\textless{}}\NormalTok{name}\OperatorTok{\textgreater{}}               \CommentTok{\# Full removal: scrubs aliases, removes config references, optionally deletes or archives}
\end{Highlighting}
\end{Shaded}

\subsection{Multiple Named Resources}\label{multiple-named-resources}

A profile can have multiple named journals and ledgers:

\begin{Shaded}
\begin{Highlighting}[]
\ExtensionTok{ww}\NormalTok{ journal add }\OperatorTok{\textless{}}\NormalTok{name}\OperatorTok{\textgreater{}}          \CommentTok{\# Creates journals/\textless{}name\textgreater{}.txt and adds mapping to jrnl.yaml}
\ExtensionTok{ww}\NormalTok{ journal list                }\CommentTok{\# Lists all journals in active profile}
\ExtensionTok{ww}\NormalTok{ journal remove }\OperatorTok{\textless{}}\NormalTok{name}\OperatorTok{\textgreater{}}       \CommentTok{\# Removes journal file and mapping}
\ExtensionTok{ww}\NormalTok{ journal rename }\OperatorTok{\textless{}}\NormalTok{old}\OperatorTok{\textgreater{}} \OperatorTok{\textless{}}\NormalTok{new}\OperatorTok{\textgreater{}}  \CommentTok{\# Renames journal file and updates mapping}

\ExtensionTok{ww}\NormalTok{ ledger add }\OperatorTok{\textless{}}\NormalTok{name}\OperatorTok{\textgreater{}}           \CommentTok{\# Creates ledgers/\textless{}name\textgreater{}.journal and adds mapping to ledgers.yaml}
\ExtensionTok{ww}\NormalTok{ ledger list                 }\CommentTok{\# Lists all ledgers in active profile}
\end{Highlighting}
\end{Shaded}

\subsection{Profile Groups}\label{profile-groups}

\begin{Shaded}
\begin{Highlighting}[]
\ExtensionTok{ww}\NormalTok{ group create }\OperatorTok{\textless{}}\NormalTok{name}\OperatorTok{\textgreater{}}                    \CommentTok{\# Create a named group}
\ExtensionTok{ww}\NormalTok{ group add }\OperatorTok{\textless{}}\NormalTok{group}\OperatorTok{\textgreater{}} \OperatorTok{\textless{}}\NormalTok{profile}\OperatorTok{\textgreater{}}            \CommentTok{\# Add profile to group}
\ExtensionTok{ww}\NormalTok{ group list                             }\CommentTok{\# List all groups}
\ExtensionTok{ww}\NormalTok{ group show }\OperatorTok{\textless{}}\NormalTok{name}\OperatorTok{\textgreater{}}                      \CommentTok{\# Show group membership}
\ExtensionTok{ww}\NormalTok{ group remove }\OperatorTok{\textless{}}\NormalTok{group}\OperatorTok{\textgreater{}} \OperatorTok{\textless{}}\NormalTok{profile}\OperatorTok{\textgreater{}}         \CommentTok{\# Remove profile from group}
\end{Highlighting}
\end{Shaded}

Groups enable batch operations: \texttt{ww\ export\ -\/-group\ work}
exports all profiles in the \texttt{work} group.

\subsection{Backup/Restore Policy}\label{backuprestore-policy}

Profile backup creates a \texttt{.tar.gz} archive of the entire profile
directory. The archive is self-contained ---
\texttt{ww\ profile\ import\ \textless{}archive\textgreater{}}
reconstructs the profile including all data, config, and resource
mappings.

Safety backups are created automatically before \texttt{profile\ delete}
and \texttt{ww\ remove}. The backup path is displayed on screen before
deletion proceeds.

\begin{center}\rule{0.5\linewidth}{0.5pt}\end{center}

\section{6. Shell Bootstrap}\label{shell-bootstrap}

\subsection{bin/ww-init.sh}\label{binww-init.sh}

Sourced by the user's \texttt{.bashrc} or \texttt{.zshrc} at shell
start. Responsibilities:

\begin{enumerate}
\def\labelenumi{\arabic{enumi}.}
\tightlist
\item
  Sources \texttt{lib/shell-integration.sh}
\item
  Calls \texttt{ensure\_shell\_functions()} to verify injection is up to
  date
\item
  Restores the last active profile if \texttt{WARRIOR\_PROFILE} is unset
\end{enumerate}

\subsection{Shell Functions Injected}\label{shell-functions-injected}

{\def\LTcaptype{none} % do not increment counter
\begin{longtable}[]{@{}
  >{\raggedright\arraybackslash}p{(\linewidth - 2\tabcolsep) * \real{0.4348}}
  >{\raggedright\arraybackslash}p{(\linewidth - 2\tabcolsep) * \real{0.5652}}@{}}
\toprule\noalign{}
\begin{minipage}[b]{\linewidth}\raggedright
Function
\end{minipage} & \begin{minipage}[b]{\linewidth}\raggedright
Description
\end{minipage} \\
\midrule\noalign{}
\endhead
\bottomrule\noalign{}
\endlastfoot
\texttt{task\ {[}args{]}} & TaskWarrior with active profile's TASKRC and
TASKDATA \\
\texttt{timew\ {[}args{]}} & TimeWarrior with active profile's
TIMEWARRIORDB \\
\texttt{j\ {[}name{]}\ "text"} & JRNL entry --- optional journal name,
falls back to default \\
\texttt{l\ {[}name{]}\ \textless{}cmd\textgreater{}} & Hledger command
--- optional ledger name, falls back to default \\
\texttt{i\ {[}args{]}} & Bugwarrior pull shorthand \\
\texttt{q\ \textless{}template\textgreater{}} & Questions template
runner \\
\texttt{list\ {[}args{]}} & Quick task list (ww list passthrough) \\
\texttt{search\ \textless{}term\textgreater{}} & Cross-profile search \\
\texttt{p-\textless{}name\textgreater{}} & Activate named profile ---
created per profile \\
\end{longtable}
}

\subsection{RC File Management}\label{rc-file-management}

\texttt{shell-integration.sh} manages a \texttt{\#\ WW\ ALIASES} section
in all detected shell rc files (\texttt{.bashrc}, \texttt{.zshrc}).
Profile aliases are written to this section.
\texttt{remove\_profile\_aliases()} removes only the entries for the
specified profile. The section is idempotent --- re-running never
duplicates entries.

\begin{center}\rule{0.5\linewidth}{0.5pt}\end{center}

\section{7. CLI Dispatcher}\label{cli-dispatcher}

\subsection{bin/ww}\label{binww}

The main entry point for all \texttt{ww} commands. Approximately 709
lines. Responsibilities:

\begin{enumerate}
\def\labelenumi{\arabic{enumi}.}
\tightlist
\item
  Sources all \texttt{lib/*.sh} files
\item
  Checks profile active state (or allow-listed command flag)
\item
  Parses first argument as the service category
\item
  Discovers service script: checks
  \texttt{\$WORKWARRIOR\_BASE/services/\textless{}category\textgreater{}/}
  first, then \texttt{services/\textless{}category\textgreater{}/}
\item
  Validates known subcommands at the routing level for
  security-sensitive services
\item
  Executes the service with remaining arguments
\end{enumerate}

\subsection{Service Discovery}\label{service-discovery}

The dispatcher scans for executables matching the service category.
Profile-level services shadow global services of the same name. If no
match is found, the dispatcher prints a list of available services and
exits 1.

\subsection{Allow-Listed Commands}\label{allow-listed-commands}

Some commands run without an active profile: \texttt{profile\ create},
\texttt{profile\ list}, \texttt{deps\ install}, \texttt{deps\ check},
\texttt{help}, \texttt{version}. All others require
\texttt{WARRIOR\_PROFILE} to be set.

\subsection{Security Constraint}\label{security-constraint}

For the browser CMD endpoint, \texttt{ALLOWED\_SUBCOMMANDS} is a
frozenset of valid service names. POST /cmd requests must have a first
token matching this set. No \texttt{sh\ -c}, no eval. Unknown
subcommands return HTTP 400.

\begin{center}\rule{0.5\linewidth}{0.5pt}\end{center}

\section{8. Core Libraries}\label{core-libraries}

\subsection{lib/core-utils.sh}\label{libcore-utils.sh}

\begin{itemize}
\tightlist
\item
  \texttt{resolve\_ww\_base()} --- resolves \texttt{WORKWARRIOR\_BASE}
  from \texttt{WARRIOR\_PROFILE} if not set
\item
  \texttt{validate\_profile\_active()} --- exits 1 with message if no
  profile is active
\item
  \texttt{get\_last\_profile()} / \texttt{set\_last\_profile()} ---
  persists last-used profile name to
  \texttt{\textasciitilde{}/.ww\_last\_profile}
\item
  \texttt{profile\_exists(name)} --- returns 0 if profile directory
  exists
\end{itemize}

\subsection{lib/profile-manager.sh}\label{libprofile-manager.sh}

\begin{itemize}
\tightlist
\item
  \texttt{create\_profile(name)} --- full profile creation: mkdir
  structure, write \texttt{.taskrc} template, inject TaskWarrior hooks,
  create \texttt{jrnl.yaml} and \texttt{ledgers.yaml}, write bugwarrior
  config template, call \texttt{create\_profile\_aliases()}
\item
  \texttt{delete\_profile(name)} --- creates safety backup, removes
  profile directory
\item
  \texttt{backup\_profile(name)} --- archives to
  \texttt{profiles/\textless{}name\textgreater{}-backup-\textless{}timestamp\textgreater{}.tar.gz}
\item
  \texttt{import\_profile(archive)} --- extracts archive to
  \texttt{profiles/}, creates aliases
\item
  \texttt{restore\_profile(archive)} --- replaces existing profile
  directory from archive
\end{itemize}

\subsection{lib/shell-integration.sh}\label{libshell-integration.sh}

See \hyperref[6-shell-bootstrap]{Section 6}.

\begin{itemize}
\tightlist
\item
  \texttt{use\_task\_profile(name)} --- core activation: exports 5 env
  vars, sets last profile
\item
  \texttt{create\_profile\_aliases(name)} --- writes
  \texttt{p-\textless{}name\textgreater{}} to all rc files
\item
  \texttt{remove\_profile\_aliases(name)} --- removes all aliases for
  profile from rc files
\item
  \texttt{get\_ww\_rc\_files()} --- returns list of shell rc files to
  write to
\item
  \texttt{add\_alias\_to\_section(name,\ value,\ {[}file{]})} ---
  idempotent alias injection
\item
  \texttt{ensure\_shell\_functions()} --- verifies \texttt{ww-init.sh}
  source line present in all rc files
\end{itemize}

Re-source guard:
\texttt{{[}{[}\ -n\ "\$\{SHELL\_INTEGRATION\_LOADED:-\}"\ {]}{]}\ \&\&\ return\ 0}.
Never \texttt{readonly} --- normal user re-sourcing must not error.

\subsection{lib/logging.sh}\label{liblogging.sh}

All user-facing output goes through these functions. Never use raw
\texttt{echo} in lib or services.

{\def\LTcaptype{none} % do not increment counter
\begin{longtable}[]{@{}lll@{}}
\toprule\noalign{}
Function & Output & Exit \\
\midrule\noalign{}
\endhead
\bottomrule\noalign{}
\endlastfoot
\texttt{log\_info\ "msg"} & \texttt{{[}INFO{]}\ msg} to stdout & --- \\
\texttt{log\_success\ "msg"} & \texttt{{[}✓{]}\ msg} to stdout & --- \\
\texttt{log\_warning\ "msg"} & \texttt{{[}WARN{]}\ msg} to stderr &
--- \\
\texttt{log\_error\ "msg"} & \texttt{{[}ERROR{]}\ msg} to stderr &
--- \\
\texttt{log\_step\ "msg"} & \texttt{{[}→{]}\ msg} to stdout & --- \\
\end{longtable}
}

\subsection{lib/taskwarrior-api.sh}\label{libtaskwarrior-api.sh}

Wrappers for TaskWarrior operations that avoid direct \texttt{task}
calls in lib functions (keeping testability):

\begin{itemize}
\tightlist
\item
  \texttt{tw\_get\_task(uuid)} --- returns task JSON
\item
  \texttt{tw\_update\_task(uuid,\ mods)} --- applies modifications
\item
  \texttt{tw\_get\_field(uuid,\ field)} --- returns single field value
\end{itemize}

\subsection{lib/config-utils.sh}\label{libconfig-utils.sh}

YAML parsing utilities for jrnl.yaml, ledgers.yaml, and ai.yaml:

\begin{itemize}
\tightlist
\item
  \texttt{yaml\_get(file,\ key)} --- reads a single scalar value
\item
  \texttt{yaml\_get\_mapping(file)} --- returns key:value pairs from a
  mapping
\item
  \texttt{yaml\_set(file,\ key,\ value)} --- writes or updates a value
\end{itemize}

\subsection{lib/config-loader.sh}\label{libconfig-loader.sh}

Loads and validates GitHub sync configuration from bugwarrior config:

\begin{itemize}
\tightlist
\item
  \texttt{init\_github\_sync\_config()} --- main entry point, locates
  config, loads, validates
\item
  \texttt{load\_github\_sync\_config(config\_path)} --- exports
  \texttt{GITHUB\_LOGIN}, \texttt{GITHUB\_USERNAME},
  \texttt{GITHUB\_TOKEN}, \texttt{GITHUB\_PROJECT\_TEMPLATE}
\item
  \texttt{validate\_github\_sync\_config()} --- checks required fields,
  returns 1 with specific error if missing
\item
  Oracle token pattern: \texttt{@oracle:eval:gh\ auth\ token} ---
  evaluated at load time via \texttt{gh\ auth\ token}
\end{itemize}

\begin{center}\rule{0.5\linewidth}{0.5pt}\end{center}

\section{9. Service Architecture}\label{service-architecture}

\subsection{Contract}\label{contract}

Every service script must:

\begin{enumerate}
\def\labelenumi{\arabic{enumi}.}
\tightlist
\item
  Start with \texttt{\#!/usr/bin/env\ bash} and
  \texttt{set\ -euo\ pipefail}
\item
  Respond to \texttt{-\/-help} / \texttt{-h} with a one-line description
  and usage example
\item
  Use exit codes: \texttt{0} success, \texttt{1} user error, \texttt{2}
  system/internal error
\item
  Log via \texttt{lib/logging.sh} functions --- never raw \texttt{echo}
  for user-facing messages
\item
  Not write to profile directories directly --- call lib functions
\item
  Be discoverable: placed at
  \texttt{services/\textless{}category\textgreater{}/} with executable
  permission
\end{enumerate}

\subsection{Minimal Service Template}\label{minimal-service-template}

\begin{Shaded}
\begin{Highlighting}[]
\CommentTok{\#!/usr/bin/env bash}
\BuiltInTok{set} \AttributeTok{{-}euo}\NormalTok{ pipefail}

\BuiltInTok{source} \StringTok{"}\VariableTok{$WORKWARRIOR\_BASE}\StringTok{/lib/logging.sh"}

\VariableTok{USAGE}\OperatorTok{=}\StringTok{"ww myservice — one{-}line description}
\StringTok{Usage: ww myservice \textless{}subcommand\textgreater{} [args]}
\StringTok{"}

\ControlFlowTok{case} \StringTok{"}\VariableTok{$\{1}\OperatorTok{:{-}}\VariableTok{\}}\StringTok{"} \KeywordTok{in}
  \SpecialStringTok{{-}{-}help}\KeywordTok{|}\SpecialStringTok{{-}h}\KeywordTok{)}  \BuiltInTok{echo} \StringTok{"}\VariableTok{$USAGE}\StringTok{"}\KeywordTok{;} \BuiltInTok{exit}\NormalTok{ 0 }\ControlFlowTok{;;}
  \SpecialStringTok{list}\KeywordTok{)}       \CommentTok{\# implementation ;;}
  \ExtensionTok{info}\ErrorTok{)}
    \KeywordTok{[[} \OtherTok{{-}z} \StringTok{"}\VariableTok{$\{2}\OperatorTok{:{-}}\VariableTok{\}}\StringTok{"} \KeywordTok{]]} \KeywordTok{\&\&} \KeywordTok{\{} \ExtensionTok{log\_error} \StringTok{"info requires a name"}\KeywordTok{;} \BuiltInTok{exit}\NormalTok{ 1}\KeywordTok{;} \KeywordTok{\}}
    \CommentTok{\# implementation ;;}
  \ExtensionTok{*}\ErrorTok{)}
    \ExtensionTok{log\_error} \StringTok{"Unknown subcommand: \textquotesingle{}}\VariableTok{$\{1}\OperatorTok{:{-}}\VariableTok{\}}\StringTok{\textquotesingle{}"}
    \BuiltInTok{echo} \StringTok{"}\VariableTok{$USAGE}\StringTok{"}
    \BuiltInTok{exit}\NormalTok{ 1}
    \ControlFlowTok{;;}
\ControlFlowTok{esac}
\end{Highlighting}
\end{Shaded}

\subsection{Profile-Level Service
Override}\label{profile-level-service-override}

Services at
\texttt{profiles/\textless{}name\textgreater{}/services/\textless{}category\textgreater{}/}
shadow global services. The dispatcher checks the profile path first.
This enables per-profile service customizations without modifying global
services.

\subsection{Service Tiers}\label{service-tiers}

{\def\LTcaptype{none} % do not increment counter
\begin{longtable}[]{@{}
  >{\raggedright\arraybackslash}p{(\linewidth - 4\tabcolsep) * \real{0.2069}}
  >{\raggedright\arraybackslash}p{(\linewidth - 4\tabcolsep) * \real{0.4483}}
  >{\raggedright\arraybackslash}p{(\linewidth - 4\tabcolsep) * \real{0.3448}}@{}}
\toprule\noalign{}
\begin{minipage}[b]{\linewidth}\raggedright
Tier
\end{minipage} & \begin{minipage}[b]{\linewidth}\raggedright
Description
\end{minipage} & \begin{minipage}[b]{\linewidth}\raggedright
Use Case
\end{minipage} \\
\midrule\noalign{}
\endhead
\bottomrule\noalign{}
\endlastfoot
1 --- Simple & Single-file, direct tool calls & Single-tool wrappers,
simple data reads \\
2 --- Compound & Multiple subcommands, uses lib functions & Lifecycle
services (create/list/delete/backup) \\
3 --- Complex & State management, external APIs, background processes &
GitHub sync, browser server \\
\end{longtable}
}

\begin{center}\rule{0.5\linewidth}{0.5pt}\end{center}

\section{10. Service Registry}\label{service-registry}

{\def\LTcaptype{none} % do not increment counter
\begin{longtable}[]{@{}
  >{\raggedright\arraybackslash}p{(\linewidth - 4\tabcolsep) * \real{0.2581}}
  >{\raggedright\arraybackslash}p{(\linewidth - 4\tabcolsep) * \real{0.3226}}
  >{\raggedright\arraybackslash}p{(\linewidth - 4\tabcolsep) * \real{0.4194}}@{}}
\toprule\noalign{}
\begin{minipage}[b]{\linewidth}\raggedright
Domain
\end{minipage} & \begin{minipage}[b]{\linewidth}\raggedright
Commands
\end{minipage} & \begin{minipage}[b]{\linewidth}\raggedright
Description
\end{minipage} \\
\midrule\noalign{}
\endhead
\bottomrule\noalign{}
\endlastfoot
\texttt{ww\ profile} & create, list, info, delete, backup, import,
restore, uda, urgency, density & Profile lifecycle and configuration \\
\texttt{ww\ journal} & add, list, remove, rename & Named journal
management \\
\texttt{ww\ ledger} & add, list, remove, rename & Named ledger
management \\
\texttt{ww\ group} & list, create, show, add, remove, delete & Profile
groups for batch operations \\
\texttt{ww\ model} & list, providers, show, add-provider,
remove-provider, add-model, set-default, env, check & LLM provider and
model registry \\
\texttt{ww\ ctrl} & status, ai-on, ai-off, ai-status, ai-mode,
prompt-ww, prompt-ai & AI mode, prompt settings, UI config \\
\texttt{ww\ find} & \texttt{\textless{}term\textgreater{}} with filters
& Cross-profile search across tasks, time, journals, ledgers \\
\texttt{ww\ issues} & sync, push, pull, status, enable, disable, custom,
uda & GitHub two-way sync + Bugwarrior pull \\
\texttt{ww\ custom} & journals, ledgers, tasks, times, issues &
Interactive configuration wizards \\
\texttt{ww\ extensions} & taskwarrior list/search/info/refresh &
Extension registry \\
\texttt{ww\ export} & JSON, CSV, markdown & Profile data export \\
\texttt{ww\ questions} & list, new, delete,
\texttt{\textless{}template\textgreater{}} & Template-based capture
workflows \\
\texttt{ww\ browser} & start, stop, status, export & Locally-served web
UI \\
\texttt{ww\ remove} & \texttt{\textless{}profile\textgreater{}}, --keep,
--all, --archive-all, --delete-all & Profile removal with scrubbing \\
\texttt{ww\ shortcut} & list, info, add, remove & Shortcut/alias
reference \\
\texttt{ww\ deps} & install, check & Dependency management \\
\texttt{ww\ compile-heuristics} & --verbose, --digest & Recompile
NL→command rules \\
\texttt{ww\ gun} & \texttt{\textless{}args\textgreater{}} & Bulk task
series (wraps taskgun) \\
\texttt{ww\ sword} & \texttt{\textless{}task\textgreater{}} -p N
{[}--interval Nd{]} & Task splitting with dependency chains \\
\texttt{ww\ next} & --- & CFS-inspired next-task recommendation \\
\texttt{ww\ schedule} & --- & Auto-scheduler (wraps taskcheck) \\
\texttt{ww\ mcp} & install, status & MCP server for AI agent access \\
\texttt{ww\ tui} & install & taskwarrior-tui installer \\
\texttt{ww\ list} & --- & List management tool \\
\end{longtable}
}

\begin{center}\rule{0.5\linewidth}{0.5pt}\end{center}

\section{11. Browser UI}\label{browser-ui}

\subsection{Overview}\label{overview-1}

\texttt{ww\ browser} starts a locally-served web interface on
\texttt{http://localhost:7777}. No npm, no build step, no cloud, no
accounts. Python 3 stdlib only.

\begin{Shaded}
\begin{Highlighting}[]
\ExtensionTok{ww}\NormalTok{ browser                     }\CommentTok{\# Start on port 7777, open browser}
\ExtensionTok{ww}\NormalTok{ browser }\AttributeTok{{-}{-}port}\NormalTok{ 9090         }\CommentTok{\# Custom port}
\ExtensionTok{ww}\NormalTok{ browser }\AttributeTok{{-}{-}no{-}open}           \CommentTok{\# Start without opening browser}
\ExtensionTok{ww}\NormalTok{ browser stop                }\CommentTok{\# Stop server}
\ExtensionTok{ww}\NormalTok{ browser status              }\CommentTok{\# Show running state}
\end{Highlighting}
\end{Shaded}

\subsection{Server Implementation}\label{server-implementation}

\begin{itemize}
\tightlist
\item
  \textbf{\texttt{services/browser/server.py}} --- single Python file
  implementing \texttt{ThreadingHTTPServer}
\item
  Static files served from \texttt{services/browser/static/} on each
  request (no caching --- changes visible on refresh)
\item
  \texttt{ThreadingHTTPServer} handles SSE connections (which hold
  sockets open) without blocking concurrent POST requests
\item
  SSE endpoint: \texttt{GET\ /events} --- streams profile change events
  as \texttt{text/event-stream}
\end{itemize}

\subsection{REST API}\label{rest-api}

{\def\LTcaptype{none} % do not increment counter
\begin{longtable}[]{@{}lll@{}}
\toprule\noalign{}
Method & Path & Description \\
\midrule\noalign{}
\endhead
\bottomrule\noalign{}
\endlastfoot
GET & /health & Liveness check with profile name and version \\
GET & /data/tasks & Pending tasks for active profile \\
GET & /data/time & Time intervals and totals \\
GET & /data/journal & Recent journal entries (active journal) \\
GET & /data/ledger & Account balances and recent transactions \\
GET & /data/ctrl & AI settings and resolved provider \\
GET & /data/models & LLM provider registry \\
GET & /data/profile & Active profile info and resource lists \\
GET & /data/groups & Profile group definitions \\
GET & /data/sync & GitHub sync state for active profile \\
GET & /data/questions & Available question templates \\
POST & /cmd & Execute a ww subcommand \\
POST & /profile/switch & Switch active profile \\
GET & /events & SSE stream \\
\end{longtable}
}

\subsection{Security Model}\label{security-model}

\texttt{ALLOWED\_SUBCOMMANDS} frozenset: every POST /cmd request is
validated --- the first token must appear in the frozenset. Unknown
subcommands return HTTP 400. No \texttt{sh\ -c}, no eval.

This bounds the attack surface: XSS in the UI can only invoke known ww
subcommands. Known ww subcommands cannot exec arbitrary shell commands.

\subsection{Panels}\label{panels}

{\def\LTcaptype{none} % do not increment counter
\begin{longtable}[]{@{}
  >{\raggedright\arraybackslash}p{(\linewidth - 4\tabcolsep) * \real{0.3182}}
  >{\raggedright\arraybackslash}p{(\linewidth - 4\tabcolsep) * \real{0.2727}}
  >{\raggedright\arraybackslash}p{(\linewidth - 4\tabcolsep) * \real{0.4091}}@{}}
\toprule\noalign{}
\begin{minipage}[b]{\linewidth}\raggedright
Panel
\end{minipage} & \begin{minipage}[b]{\linewidth}\raggedright
Data
\end{minipage} & \begin{minipage}[b]{\linewidth}\raggedright
Actions
\end{minipage} \\
\midrule\noalign{}
\endhead
\bottomrule\noalign{}
\endlastfoot
Tasks & Full task list with UDAs & Inline edit, start/stop/done, add,
annotate \\
Time & Today's total, weekly breakdown, recent intervals & Start,
stop \\
Journals & Entry list with expand/collapse & New entry, journal selector
dropdown \\
Ledgers & Account balances, recent transactions, income statement & New
transaction, ledger selector \\
CMD & Natural language + direct command input & Route indicator ⚡/⚙ \\
CTRL & AI mode toggle & off / local-only / local+remote \\
Models & LLM provider and model registry & Add provider, set default \\
Groups & Profile group management & Create, add member \\
Sync & GitHub sync dashboard & Sync, push, status \\
Questions & Template browser & Run template \\
Profile & Profile info and resource management & Switch profile,
resource list \\
Weapons bar & Gun, Sword, Next, Schedule icons & Opens weapon panel \\
\end{longtable}
}

\subsection{CMD Input Routing}\label{cmd-input-routing}

\begin{verbatim}
Input received
  → Compound check: contains "and"/"then"/"also"/"plus"?
    → Yes: split into segments
  → Each segment:
    → HeuristicEngine.match(segment, threshold=0.8)
    → If match: execute
    → If no match: AI fallback if configured
    → If no AI: return error with suggestion
  → Collect results, return to UI
\end{verbatim}

\begin{center}\rule{0.5\linewidth}{0.5pt}\end{center}

\section{12. Heuristic Engine}\label{heuristic-engine}

\subsection{Overview}\label{overview-2}

627 compiled regex rules across 19 command domains. No network. No
latency. Loaded at server startup from
\texttt{config/cmd-heuristics.yaml}.

\subsection{Pipeline}\label{pipeline}

\begin{enumerate}
\def\labelenumi{\arabic{enumi}.}
\tightlist
\item
  \textbf{Compound split} --- input containing ``and'', ``then'',
  ``also'', ``plus'' is split into segments
\item
  \textbf{Rule evaluation} --- each segment tested against all 627
  rules, highest confidence wins
\item
  \textbf{Threshold check} --- match must exceed 0.8 confidence to
  accept
\item
  \textbf{AI fallback} --- if no rule matches, input goes to configured
  LLM (if enabled)
\item
  \textbf{Execution} --- matched commands executed against active
  profile
\end{enumerate}

\subsection{Rule Structure}\label{rule-structure}

Each rule contains: - \texttt{pattern} --- compiled regex -
\texttt{action\_template} --- output command template with capture group
references - \texttt{confidence} --- score (0.0--1.0) - \texttt{domain}
--- one of 19 command domains

\subsection{Confidence Levels by Phrasing
Variation}\label{confidence-levels-by-phrasing-variation}

{\def\LTcaptype{none} % do not increment counter
\begin{longtable}[]{@{}
  >{\raggedright\arraybackslash}p{(\linewidth - 4\tabcolsep) * \real{0.3548}}
  >{\raggedright\arraybackslash}p{(\linewidth - 4\tabcolsep) * \real{0.3548}}
  >{\raggedright\arraybackslash}p{(\linewidth - 4\tabcolsep) * \real{0.2903}}@{}}
\toprule\noalign{}
\begin{minipage}[b]{\linewidth}\raggedright
Variation
\end{minipage} & \begin{minipage}[b]{\linewidth}\raggedright
Confidence
\end{minipage} & \begin{minipage}[b]{\linewidth}\raggedright
Example
\end{minipage} \\
\midrule\noalign{}
\endhead
\bottomrule\noalign{}
\endlastfoot
Passthrough & 1.0 & \texttt{task\ add\ review\ budget} \\
Imperative & 0.95 & \texttt{add\ a\ task\ to\ review\ the\ budget} \\
Declarative & 0.90 &
\texttt{I\ need\ a\ task\ for\ reviewing\ the\ budget} \\
Interrogative & 0.90 &
\texttt{can\ you\ create\ a\ task\ to\ review\ the\ budget} \\
Shorthand & 0.90 & \texttt{task:\ review\ budget\ due\ friday} \\
Verbose & 0.85 &
\texttt{I\ would\ like\ to\ add\ a\ new\ task\ for\ reviewing\ the\ budget} \\
\end{longtable}
}

\subsection{Command Domains}\label{command-domains}

task · time · journal · ledger · profile · group · model · ctrl ·
service · issues · find · schedule · gun · next · mcp · browser ·
extensions · custom · shortcut

\subsection{Date Expression
Normalization}\label{date-expression-normalization}

{\def\LTcaptype{none} % do not increment counter
\begin{longtable}[]{@{}ll@{}}
\toprule\noalign{}
Input & Output \\
\midrule\noalign{}
\endhead
\bottomrule\noalign{}
\endlastfoot
``tomorrow'' & \texttt{due:tomorrow} \\
``next week'' & \texttt{due:1w} \\
``friday'' & \texttt{due:friday} \\
``end of month'' & \texttt{due:eom} \\
``in 3 days'' & \texttt{due:3d} \\
\end{longtable}
}

\subsection{Compiler}\label{compiler}

\texttt{scripts/compile-heuristics.py} generates
\texttt{config/cmd-heuristics.yaml}:

\begin{Shaded}
\begin{Highlighting}[]
\ExtensionTok{ww}\NormalTok{ compile{-}heuristics              }\CommentTok{\# Standard run}
\ExtensionTok{ww}\NormalTok{ compile{-}heuristics }\AttributeTok{{-}{-}verbose}    \CommentTok{\# Every rule with test results}
\ExtensionTok{ww}\NormalTok{ compile{-}heuristics }\AttributeTok{{-}{-}digest}     \CommentTok{\# + CMD log analysis → new rules from AI hits}
\end{Highlighting}
\end{Shaded}

Compiler reads: \texttt{bin/ww} case branches,
\texttt{config/shortcuts.yaml}, \texttt{config/command-syntax.yaml}, and
optionally \texttt{services/cmd/cmd.log}.

Steps: 1. Scan command sources for command definitions 2. Generate regex
patterns for each phrasing variation 3. Validate against synthetic
corpus (\texttt{config/cmd-heuristics-corpus.yaml}) 4. Resolve conflicts
(higher-confidence rule wins) 5. Fill coverage gaps 6. Write output

\subsection{Self-Improvement Loop}\label{self-improvement-loop}

Every CMD submission is logged to \texttt{services/cmd/cmd.log} as
JSONL: \texttt{\{input,\ route,\ output,\ success,\ timestamp\}}.

\texttt{-\/-digest} flag reads this log, finds successful AI
translations that had no heuristic match, generates new rules from those
translations. Run \texttt{ww\ compile-heuristics\ -\/-digest} regularly
to expand heuristic coverage from real usage.

\begin{center}\rule{0.5\linewidth}{0.5pt}\end{center}

\section{13. AI Integration}\label{ai-integration}

\subsection{Configuration}\label{configuration}

\begin{Shaded}
\begin{Highlighting}[]
\CommentTok{\# config/ai.yaml}
\FunctionTok{mode}\KeywordTok{:}\AttributeTok{ }\CharTok{off}\CommentTok{                    \# off | local{-}only | local+remote}
\FunctionTok{preferred\_provider}\KeywordTok{:}\AttributeTok{ ollama}
\FunctionTok{access\_points}\KeywordTok{:}
\AttributeTok{  }\FunctionTok{cmd\_ai}\KeywordTok{:}\AttributeTok{ }\CharTok{true}\CommentTok{               \# Enable AI in CMD service}
\end{Highlighting}
\end{Shaded}

Per-profile override:
\texttt{profiles/\textless{}name\textgreater{}/ai.yaml} with same
structure. Profile config takes precedence over global.

\subsection{Modes}\label{modes}

{\def\LTcaptype{none} % do not increment counter
\begin{longtable}[]{@{}
  >{\raggedright\arraybackslash}p{(\linewidth - 2\tabcolsep) * \real{0.4000}}
  >{\raggedright\arraybackslash}p{(\linewidth - 2\tabcolsep) * \real{0.6000}}@{}}
\toprule\noalign{}
\begin{minipage}[b]{\linewidth}\raggedright
Mode
\end{minipage} & \begin{minipage}[b]{\linewidth}\raggedright
Behavior
\end{minipage} \\
\midrule\noalign{}
\endhead
\bottomrule\noalign{}
\endlastfoot
\texttt{off} & Heuristics only. No AI calls. Unmatched inputs return
error. \\
\texttt{local-only} & Heuristics first, then ollama if no match. No data
leaves machine. \\
\texttt{local+remote} & Heuristics → local LLM → remote provider
fallback chain. \\
\end{longtable}
}

\subsection{Provider Registry}\label{provider-registry}

\begin{Shaded}
\begin{Highlighting}[]
\CommentTok{\# config/models.yaml}
\FunctionTok{providers}\KeywordTok{:}
\AttributeTok{  }\KeywordTok{{-}}\AttributeTok{ }\FunctionTok{name}\KeywordTok{:}\AttributeTok{ ollama}
\AttributeTok{    }\FunctionTok{type}\KeywordTok{:}\AttributeTok{ ollama}
\AttributeTok{    }\FunctionTok{base\_url}\KeywordTok{:}\AttributeTok{ http://localhost:11434}
\AttributeTok{  }\KeywordTok{{-}}\AttributeTok{ }\FunctionTok{name}\KeywordTok{:}\AttributeTok{ openai}
\AttributeTok{    }\FunctionTok{type}\KeywordTok{:}\AttributeTok{ openai}
\AttributeTok{    }\FunctionTok{api\_key\_env}\KeywordTok{:}\AttributeTok{ OPENAI\_API\_KEY}
\FunctionTok{models}\KeywordTok{:}
\AttributeTok{  }\FunctionTok{default}\KeywordTok{:}\AttributeTok{ llama3.2}
\AttributeTok{  }\FunctionTok{fallback\_chain}\KeywordTok{:}\AttributeTok{ }\KeywordTok{[}\AttributeTok{llama3}\FloatTok{.2}\KeywordTok{,}\AttributeTok{ gpt{-}4o{-}mini}\KeywordTok{]}
\end{Highlighting}
\end{Shaded}

CLI management:

\begin{Shaded}
\begin{Highlighting}[]
\ExtensionTok{ww}\NormalTok{ model add{-}provider ollama ollama http://localhost:11434}
\ExtensionTok{ww}\NormalTok{ model add{-}provider openai openai}
\ExtensionTok{ww}\NormalTok{ model set{-}default llama3.2}
\ExtensionTok{ww}\NormalTok{ model list}
\ExtensionTok{ww}\NormalTok{ model check          }\CommentTok{\# Test connectivity to all configured providers}
\end{Highlighting}
\end{Shaded}

\subsection{Runtime Control}\label{runtime-control}

\begin{Shaded}
\begin{Highlighting}[]
\ExtensionTok{ww}\NormalTok{ ctrl ai{-}on           }\CommentTok{\# Enable (uses mode from config)}
\ExtensionTok{ww}\NormalTok{ ctrl ai{-}off          }\CommentTok{\# Disable}
\ExtensionTok{ww}\NormalTok{ ctrl ai{-}status       }\CommentTok{\# Show resolved state (global + profile + effective)}
\ExtensionTok{ww}\NormalTok{ ctrl ai{-}mode local{-}only  }\CommentTok{\# Set mode}
\end{Highlighting}
\end{Shaded}

Browser CTRL panel mirrors all CLI controls. Changes take effect
immediately --- no restart required.

\begin{center}\rule{0.5\linewidth}{0.5pt}\end{center}

\section{14. GitHub Sync Engine}\label{github-sync-engine}

\subsection{Two Engines}\label{two-engines}

\textbf{Bugwarrior} --- one-way pull from 20+ services into
TaskWarrior.\\
\textbf{ww github-sync} --- two-way bidirectional sync between
individual tasks and GitHub issues.

These are complementary. Use both simultaneously.

\subsection{Bugwarrior}\label{bugwarrior}

Pulls issues from: GitHub, GitLab, Jira, Trello, Bitbucket, Taiga,
Pagure, and 13+ more.

\begin{Shaded}
\begin{Highlighting}[]
\ExtensionTok{ww}\NormalTok{ issues custom          }\CommentTok{\# Interactive configuration wizard (per{-}profile)}
\ExtensionTok{i}\NormalTok{ pull                    }\CommentTok{\# Pull from all configured services}
\ExtensionTok{i}\NormalTok{ status                  }\CommentTok{\# Show sync state}
\end{Highlighting}
\end{Shaded}

Configuration stored in
\texttt{profiles/\textless{}name\textgreater{}/.config/bugwarrior/bugwarriorrc}.

Injected UDAs on created tasks: \texttt{githubissue},
\texttt{githuburl}, \texttt{githubrepo}, \texttt{githubauthor} ---
classified as \texttt{{[}github{]}} source in the UDA registry.

\subsection{ww github-sync}\label{ww-github-sync}

Two-way sync between individual TaskWarrior tasks and GitHub issues.

\begin{Shaded}
\begin{Highlighting}[]
\ExtensionTok{ww}\NormalTok{ issues enable }\OperatorTok{\textless{}}\NormalTok{task\_id}\OperatorTok{\textgreater{}} \OperatorTok{\textless{}}\NormalTok{issue\_num}\OperatorTok{\textgreater{}} \OperatorTok{\textless{}}\NormalTok{org/repo}\OperatorTok{\textgreater{}}    \CommentTok{\# Link task to issue}
\ExtensionTok{ww}\NormalTok{ issues disable }\OperatorTok{\textless{}}\NormalTok{task\_id}\OperatorTok{\textgreater{}}                          \CommentTok{\# Unlink}
\ExtensionTok{ww}\NormalTok{ issues sync                                        }\CommentTok{\# Two{-}way sync all linked tasks}
\ExtensionTok{ww}\NormalTok{ issues push                                        }\CommentTok{\# Push local changes to GitHub only}
\ExtensionTok{ww}\NormalTok{ issues pull                                        }\CommentTok{\# Pull GitHub changes only}
\ExtensionTok{ww}\NormalTok{ issues status                                      }\CommentTok{\# Show sync state for all linked tasks}
\end{Highlighting}
\end{Shaded}

\subsection{Field Mapping}\label{field-mapping}

\textbf{Bidirectional (both directions):}

{\def\LTcaptype{none} % do not increment counter
\begin{longtable}[]{@{}
  >{\raggedright\arraybackslash}p{(\linewidth - 2\tabcolsep) * \real{0.5000}}
  >{\raggedright\arraybackslash}p{(\linewidth - 2\tabcolsep) * \real{0.5000}}@{}}
\toprule\noalign{}
\begin{minipage}[b]{\linewidth}\raggedright
TaskWarrior
\end{minipage} & \begin{minipage}[b]{\linewidth}\raggedright
GitHub Issue
\end{minipage} \\
\midrule\noalign{}
\endhead
\bottomrule\noalign{}
\endlastfoot
\texttt{description} & title (truncated to 256 chars on push) \\
\texttt{status} & state (pending/started → OPEN, completed/deleted →
CLOSED) \\
\texttt{priority} & labels (H → priority:high, M → priority:medium, L →
priority:low) \\
\texttt{tags} & labels (system tags excluded) \\
\texttt{annotations} & comments (with \texttt{{[}tw{]}} prefix to
prevent loop) \\
\end{longtable}
}

\textbf{Pull-only (GitHub → TaskWarrior):}

{\def\LTcaptype{none} % do not increment counter
\begin{longtable}[]{@{}ll@{}}
\toprule\noalign{}
GitHub Issue & TaskWarrior UDA \\
\midrule\noalign{}
\endhead
\bottomrule\noalign{}
\endlastfoot
Issue number & \texttt{githubissue} \\
Issue URL & \texttt{githuburl} \\
Repository & \texttt{githubrepo} \\
Author & \texttt{githubauthor} \\
\texttt{createdAt} & \texttt{entry} (first sync only) \\
\texttt{closedAt} & \texttt{end} (if status → completed) \\
\texttt{updatedAt} & \texttt{modified} \\
\end{longtable}
}

\subsection{Sync Engine Files}\label{sync-engine-files}

All 10 files classified HIGH FRAGILITY. Changes require extended risk
brief.

{\def\LTcaptype{none} % do not increment counter
\begin{longtable}[]{@{}
  >{\raggedright\arraybackslash}p{(\linewidth - 2\tabcolsep) * \real{0.5000}}
  >{\raggedright\arraybackslash}p{(\linewidth - 2\tabcolsep) * \real{0.5000}}@{}}
\toprule\noalign{}
\begin{minipage}[b]{\linewidth}\raggedright
File
\end{minipage} & \begin{minipage}[b]{\linewidth}\raggedright
Role
\end{minipage} \\
\midrule\noalign{}
\endhead
\bottomrule\noalign{}
\endlastfoot
\texttt{lib/github-api.sh} & GitHub REST API via gh CLI:
get/update/label/comment \\
\texttt{lib/github-sync-state.sh} & Sync state persistence --- JSON per
task in \texttt{.config/ww-sync/} \\
\texttt{lib/sync-detector.sh} & Change detection --- compares current
state to last sync state \\
\texttt{lib/sync-pull.sh} & GitHub → TaskWarrior pull with orphan
detection \\
\texttt{lib/sync-push.sh} & TaskWarrior → GitHub push with label
management \\
\texttt{lib/sync-bidirectional.sh} & Orchestrates pull+push with
conflict window management \\
\texttt{lib/field-mapper.sh} & Field mapping between TW and GitHub
formats \\
\texttt{lib/conflict-resolver.sh} & Last-write-wins with configurable
window \\
\texttt{lib/annotation-sync.sh} & TW annotations ↔ GitHub comments \\
\texttt{services/custom/github-sync.sh} & Service entry point and CLI
dispatch \\
\end{longtable}
}

\subsection{Conflict Resolution}\label{conflict-resolution}

Last-write-wins with a configurable conflict window (default: 60
seconds).

If both sides changed within the conflict window, the sync reports the
conflict rather than silently overwriting. The conflict is surfaced in
\texttt{ww\ issues\ status} with both values and timestamps.

Outside the conflict window: the more recently modified side wins.

\subsection{Orphan Detection}\label{orphan-detection}

If a linked GitHub issue is deleted, \texttt{sync\_pull\_issue()}
receives \texttt{{[}not-found{]}} from the API. Rather than
hard-erroring, it logs a warning and skips the task. The user is
prompted to run
\texttt{ww\ issues\ disable\ \textless{}uuid\textgreater{}} to clean up
the link.

\subsection{Annotation/Comment Sync}\label{annotationcomment-sync}

TaskWarrior annotations are pushed as GitHub comments prefixed with
\texttt{{[}tw{]}}. GitHub comments prefixed with \texttt{{[}tw{]}} are
skipped on pull (preventing sync loops). Non-prefixed GitHub comments
are pulled as TaskWarrior annotations prefixed with
\texttt{{[}github{]}}.

\subsection{Oracle Token Pattern}\label{oracle-token-pattern}

Bugwarrior config stores:
\texttt{github.token\ =\ @oracle:eval:gh\ auth\ token}

The oracle directive is evaluated at load time via
\texttt{gh\ auth\ token}. The actual token is never written to the
config file. \texttt{lib/config-loader.sh} detects and evaluates this
pattern.

\begin{center}\rule{0.5\linewidth}{0.5pt}\end{center}

\section{15. Weapons System}\label{weapons-system}

Weapons are service-level tools that manipulate profile data in special
ways --- creating, slicing, and packaging tasks. They appear as a
weapons bar in the browser sidebar.

\subsection{Sword}\label{sword}

Native to ww. Splits a single task into N sequential subtasks with
dependency chains.

\begin{Shaded}
\begin{Highlighting}[]
\ExtensionTok{ww}\NormalTok{ sword }\OperatorTok{\textless{}}\NormalTok{task\_id}\OperatorTok{\textgreater{}}\NormalTok{ {-}p }\OperatorTok{\textless{}}\NormalTok{N}\OperatorTok{\textgreater{}}                   \CommentTok{\# Split into N parts}
\ExtensionTok{ww}\NormalTok{ sword }\OperatorTok{\textless{}}\NormalTok{task\_id}\OperatorTok{\textgreater{}}\NormalTok{ {-}p }\OperatorTok{\textless{}}\NormalTok{N}\OperatorTok{\textgreater{}}\NormalTok{ {-}{-}interval }\OperatorTok{\textless{}}\NormalTok{Nd}\OperatorTok{\textgreater{}}   \CommentTok{\# N{-}day intervals between parts}
\ExtensionTok{ww}\NormalTok{ sword }\OperatorTok{\textless{}}\NormalTok{task\_id}\OperatorTok{\textgreater{}}\NormalTok{ {-}p }\OperatorTok{\textless{}}\NormalTok{N}\OperatorTok{\textgreater{}}\NormalTok{ {-}{-}prefix }\StringTok{"Phase"}  \CommentTok{\# Custom prefix}
\end{Highlighting}
\end{Shaded}

Each subtask receives: - Description: ``Part N of: \{original
description\}'' - Parent task's project and tags - Due date: offset by N
× interval from now - Dependency: on previous subtask in the chain

Parent task is archived after split. The chain is strictly sequential
--- a subtask cannot be completed until all its dependencies are
complete.

\subsection{Gun}\label{gun}

Wraps \href{https://github.com/hamzamohdzubair/taskgun}{taskgun} (Rust).
Bulk task series generator.

\begin{Shaded}
\begin{Highlighting}[]
\ExtensionTok{ww}\NormalTok{ gun }\OperatorTok{\textless{}}\NormalTok{args}\OperatorTok{\textgreater{}}   \CommentTok{\# Arguments passed through to taskgun}
\end{Highlighting}
\end{Shaded}

Respects active profile (reads TASKRC/TASKDATA from env). Returns
taskgun output with error handling.

\subsection{Next}\label{next}

Wraps the \texttt{next} binary. CFS-inspired next-task recommendation.

\begin{Shaded}
\begin{Highlighting}[]
\ExtensionTok{ww}\NormalTok{ next         }\CommentTok{\# Recommends single optimal next task}
\end{Highlighting}
\end{Shaded}

Factors: TaskWarrior urgency score, deadline proximity, context signals.

\subsection{Schedule}\label{schedule}

Wraps \href{https://github.com/taskcheck/taskcheck}{taskcheck}.
Auto-scheduler for time blocks.

\begin{Shaded}
\begin{Highlighting}[]
\ExtensionTok{ww}\NormalTok{ schedule}
\end{Highlighting}
\end{Shaded}

\subsection{Planned Weapons}\label{planned-weapons}

{\def\LTcaptype{none} % do not increment counter
\begin{longtable}[]{@{}ll@{}}
\toprule\noalign{}
Weapon & Status \\
\midrule\noalign{}
\endhead
\bottomrule\noalign{}
\endlastfoot
Bat & Planned \\
Fire & Planned \\
Slingshot & Planned \\
\end{longtable}
}

\subsection{Weapon Constraints}\label{weapon-constraints}

All weapons: - Read TASKRC/TASKDATA from environment (profile isolation
respected) - Work via \texttt{POST\ /cmd} in the browser UI - Never
modify data outside active profile scope - Respond to \texttt{-\/-help}

\begin{center}\rule{0.5\linewidth}{0.5pt}\end{center}

\section{16. UDA System}\label{uda-system}

\subsection{Overview}\label{overview-3}

TaskWarrior User Defined Attributes are first-class in Workwarrior. A
profile can carry 100+ UDAs. The browser UI renders all UDAs in the task
inline editor dynamically.

\subsection{UDA Types}\label{uda-types}

{\def\LTcaptype{none} % do not increment counter
\begin{longtable}[]{@{}lll@{}}
\toprule\noalign{}
Type & TaskWarrior declaration & Use \\
\midrule\noalign{}
\endhead
\bottomrule\noalign{}
\endlastfoot
string & \texttt{uda.\textless{}name\textgreater{}.type=string} & Text
values \\
numeric & \texttt{uda.\textless{}name\textgreater{}.type=numeric} &
Numbers, billable hours \\
date & \texttt{uda.\textless{}name\textgreater{}.type=date} & Dates \\
duration & \texttt{uda.\textless{}name\textgreater{}.type=duration} &
Estimated time \\
\end{longtable}
}

\subsection{Source Classification}\label{source-classification}

{\def\LTcaptype{none} % do not increment counter
\begin{longtable}[]{@{}ll@{}}
\toprule\noalign{}
Badge & Source \\
\midrule\noalign{}
\endhead
\bottomrule\noalign{}
\endlastfoot
\texttt{{[}github{]}} & Injected by Bugwarrior from issue tracking
services \\
\texttt{{[}extension:\textless{}name\textgreater{}{]}} & Added by
TaskWarrior extensions or hooks \\
\texttt{{[}custom{]}} & Defined by user via
\texttt{ww\ profile\ uda\ add} \\
\end{longtable}
}

Extension-classified UDAs are protected from accidental deletion.

{\def\LTcaptype{none} % do not increment counter
\begin{longtable}[]{@{}lll@{}}
\toprule\noalign{}
Extension & UDAs & Badge \\
\midrule\noalign{}
\endhead
\bottomrule\noalign{}
\endlastfoot
TWDensity & \texttt{density}, \texttt{densitywindow} &
\texttt{{[}extension:twdensity{]}} \\
taskcheck & \texttt{estimated}, \texttt{time\_map} &
\texttt{{[}extension:taskcheck{]}} \\
\end{longtable}
}

\subsection{Management}\label{management}

\begin{Shaded}
\begin{Highlighting}[]
\ExtensionTok{ww}\NormalTok{ profile uda list                         }\CommentTok{\# All UDAs with source badges}
\ExtensionTok{ww}\NormalTok{ profile uda add }\OperatorTok{\textless{}}\NormalTok{name}\OperatorTok{\textgreater{}}                   \CommentTok{\# Interactive UDA creation wizard}
\ExtensionTok{ww}\NormalTok{ profile uda remove }\OperatorTok{\textless{}}\NormalTok{name}\OperatorTok{\textgreater{}}               \CommentTok{\# Remove (blocked for extension UDAs)}
\ExtensionTok{ww}\NormalTok{ profile uda group }\OperatorTok{\textless{}}\NormalTok{name}\OperatorTok{\textgreater{}}                \CommentTok{\# Apply UDA template group}
\ExtensionTok{ww}\NormalTok{ profile uda perm }\OperatorTok{\textless{}}\NormalTok{name}\OperatorTok{\textgreater{}}\NormalTok{ nosync          }\CommentTok{\# Set sync permission token}
\ExtensionTok{ww}\NormalTok{ profile uda perm }\OperatorTok{\textless{}}\NormalTok{name}\OperatorTok{\textgreater{}}\NormalTok{ private         }\CommentTok{\# Mark as private (no AI access)}
\ExtensionTok{ww}\NormalTok{ profile uda perm }\OperatorTok{\textless{}}\NormalTok{name}\OperatorTok{\textgreater{}}\NormalTok{ noai            }\CommentTok{\# Exclude from AI context}
\end{Highlighting}
\end{Shaded}

\subsection{Sync Permissions}\label{sync-permissions}

Permission tokens stored in
\texttt{profiles/\textless{}name\textgreater{}/.config/ww-sync-permissions.yaml}:

{\def\LTcaptype{none} % do not increment counter
\begin{longtable}[]{@{}ll@{}}
\toprule\noalign{}
Token & Effect \\
\midrule\noalign{}
\endhead
\bottomrule\noalign{}
\endlastfoot
\texttt{nosync} & Excluded from github-sync push \\
\texttt{private} & Excluded from AI context \\
\texttt{noai} & Excluded from AI context \\
\texttt{readonly} & Pull-only, never pushed \\
\end{longtable}
}

\subsection{Urgency Tuning}\label{urgency-tuning}

\begin{Shaded}
\begin{Highlighting}[]
\ExtensionTok{ww}\NormalTok{ profile urgency    }\CommentTok{\# Interactive urgency coefficient tuner}
\end{Highlighting}
\end{Shaded}

Tunable coefficients: \texttt{urgency.due.coefficient},
\texttt{urgency.priority.coefficient}, \texttt{urgency.age.coefficient},
\texttt{urgency.project.coefficient}, \texttt{urgency.tags.coefficient},
per-UDA urgency weights.

Results written to profile's \texttt{.taskrc}. Per-profile tuning means
different contexts can have different priority models.

\subsection{Density Scoring}\label{density-scoring}

\begin{Shaded}
\begin{Highlighting}[]
\ExtensionTok{ww}\NormalTok{ profile density    }\CommentTok{\# TWDensity integration — due{-}date density scoring}
\end{Highlighting}
\end{Shaded}

TWDensity extension injects \texttt{density} and \texttt{densitywindow}
UDAs. Higher density = more tasks due in the window = higher urgency.

\begin{center}\rule{0.5\linewidth}{0.5pt}\end{center}

\section{17. Extensions}\label{extensions}

\subsection{TaskWarrior Extensions
Registry}\label{taskwarrior-extensions-registry}

\begin{Shaded}
\begin{Highlighting}[]
\ExtensionTok{ww}\NormalTok{ extensions taskwarrior list              }\CommentTok{\# Browse community extensions}
\ExtensionTok{ww}\NormalTok{ extensions taskwarrior search }\OperatorTok{\textless{}}\NormalTok{term}\OperatorTok{\textgreater{}}     \CommentTok{\# Search by name or tag}
\ExtensionTok{ww}\NormalTok{ extensions taskwarrior info }\OperatorTok{\textless{}}\NormalTok{name}\OperatorTok{\textgreater{}}       \CommentTok{\# Show details}
\ExtensionTok{ww}\NormalTok{ extensions taskwarrior refresh           }\CommentTok{\# Update registry from upstream}
\end{Highlighting}
\end{Shaded}

Registry populated by \texttt{scripts/scan-taskwarrior-extensions.py}
--- scans the TaskWarrior community extension list and indexes metadata.

\subsection{Hook Events}\label{hook-events}

{\def\LTcaptype{none} % do not increment counter
\begin{longtable}[]{@{}ll@{}}
\toprule\noalign{}
Event & When \\
\midrule\noalign{}
\endhead
\bottomrule\noalign{}
\endlastfoot
\texttt{on-launch} & Every \texttt{task} invocation \\
\texttt{on-add} & New task creation \\
\texttt{on-modify} & Task field change \\
\texttt{on-exit} & After \texttt{task} command exits \\
\end{longtable}
}

Workwarrior's own integration uses \texttt{on-modify} to auto-start
TimeWarrior when \texttt{task\ start} is called.

\subsection{TimeWarrior Extensions}\label{timewarrior-extensions}

Python scripts that run against TimeWarrior data. Custom reports,
summaries, integrations. Installed to
\texttt{\textasciitilde{}/.timewarrior/extensions/}.

\begin{center}\rule{0.5\linewidth}{0.5pt}\end{center}

\section{18. Questions System}\label{questions-system}

Template-based structured capture workflows. A question template is a
YAML file defining a series of prompts that write to tasks, journal, or
ledger.

\begin{Shaded}
\begin{Highlighting}[]
\ExtensionTok{q} \OperatorTok{\textless{}}\NormalTok{template}\OperatorTok{\textgreater{}}            \CommentTok{\# Run a template (interactive prompts)}
\ExtensionTok{ww}\NormalTok{ questions list       }\CommentTok{\# List available templates}
\ExtensionTok{ww}\NormalTok{ questions new        }\CommentTok{\# Create a new template interactively}
\ExtensionTok{ww}\NormalTok{ questions delete     }\CommentTok{\# Remove a template}
\end{Highlighting}
\end{Shaded}

Templates in \texttt{services/questions/templates/}. Per-profile
templates in \texttt{profiles/\textless{}name\textgreater{}/templates/}
shadow global templates.

\subsection{Template Structure}\label{template-structure}

\begin{Shaded}
\begin{Highlighting}[]
\FunctionTok{name}\KeywordTok{:}\AttributeTok{ standup}
\FunctionTok{description}\KeywordTok{:}\AttributeTok{ Daily standup notes}
\FunctionTok{prompts}\KeywordTok{:}
\AttributeTok{  }\KeywordTok{{-}}\AttributeTok{ }\FunctionTok{field}\KeywordTok{:}\AttributeTok{ yesterday}
\AttributeTok{    }\FunctionTok{label}\KeywordTok{:}\AttributeTok{ }\StringTok{"What did you do yesterday?"}
\AttributeTok{    }\FunctionTok{target}\KeywordTok{:}\AttributeTok{ journal}
\AttributeTok{    }\FunctionTok{journal}\KeywordTok{:}\AttributeTok{ standup}
\AttributeTok{  }\KeywordTok{{-}}\AttributeTok{ }\FunctionTok{field}\KeywordTok{:}\AttributeTok{ today}
\AttributeTok{    }\FunctionTok{label}\KeywordTok{:}\AttributeTok{ }\StringTok{"What will you do today?"}
\AttributeTok{    }\FunctionTok{target}\KeywordTok{:}\AttributeTok{ task}
\AttributeTok{    }\FunctionTok{project}\KeywordTok{:}\AttributeTok{ standup}
\AttributeTok{  }\KeywordTok{{-}}\AttributeTok{ }\FunctionTok{field}\KeywordTok{:}\AttributeTok{ blockers}
\AttributeTok{    }\FunctionTok{label}\KeywordTok{:}\AttributeTok{ }\StringTok{"Any blockers?"}
\AttributeTok{    }\FunctionTok{target}\KeywordTok{:}\AttributeTok{ journal}
\AttributeTok{    }\FunctionTok{journal}\KeywordTok{:}\AttributeTok{ standup}
\end{Highlighting}
\end{Shaded}

\begin{center}\rule{0.5\linewidth}{0.5pt}\end{center}

\section{19. Export System}\label{export-system}

\begin{Shaded}
\begin{Highlighting}[]
\ExtensionTok{ww}\NormalTok{ export json          }\CommentTok{\# Full profile export as JSON}
\ExtensionTok{ww}\NormalTok{ export csv           }\CommentTok{\# Tasks as CSV}
\ExtensionTok{ww}\NormalTok{ export markdown      }\CommentTok{\# Tasks as Markdown table}
\ExtensionTok{ww}\NormalTok{ export }\AttributeTok{{-}{-}group} \OperatorTok{\textless{}}\NormalTok{g}\OperatorTok{\textgreater{}}   \CommentTok{\# Export all profiles in group}
\end{Highlighting}
\end{Shaded}

Export scopes: tasks only, time only, journals only, ledgers only, or
all resources.

Implemented in \texttt{lib/export-utils.sh}. Output to stdout by
default; \texttt{-\/-output\ \textless{}path\textgreater{}} writes to
file.

\begin{center}\rule{0.5\linewidth}{0.5pt}\end{center}

\section{20. Fragility Register}\label{fragility-register}

{\def\LTcaptype{none} % do not increment counter
\begin{longtable}[]{@{}
  >{\raggedright\arraybackslash}p{(\linewidth - 4\tabcolsep) * \real{0.4444}}
  >{\raggedright\arraybackslash}p{(\linewidth - 4\tabcolsep) * \real{0.1944}}
  >{\raggedright\arraybackslash}p{(\linewidth - 4\tabcolsep) * \real{0.3611}}@{}}
\toprule\noalign{}
\begin{minipage}[b]{\linewidth}\raggedright
Classification
\end{minipage} & \begin{minipage}[b]{\linewidth}\raggedright
Files
\end{minipage} & \begin{minipage}[b]{\linewidth}\raggedright
Constraints
\end{minipage} \\
\midrule\noalign{}
\endhead
\bottomrule\noalign{}
\endlastfoot
HIGH FRAGILITY & \texttt{lib/github-api.sh},
\texttt{lib/github-sync-state.sh}, \texttt{lib/sync-pull.sh},
\texttt{lib/sync-push.sh}, \texttt{lib/sync-bidirectional.sh},
\texttt{lib/field-mapper.sh}, \texttt{lib/sync-detector.sh},
\texttt{lib/conflict-resolver.sh}, \texttt{lib/annotation-sync.sh},
\texttt{services/custom/github-sync.sh} & Extended risk brief required.
Verifier sign-off before merge. Integration tests required. \\
SERIALIZED & \texttt{bin/ww}, \texttt{lib/shell-integration.sh} & One
writer at a time. Changes affect all users immediately. \\
NEVER COMMIT & \texttt{profiles/*/.task/taskchampion.sqlite3},
\texttt{profiles/*/.config/}, \texttt{profiles/*/list/},
\texttt{*.sqlite3} & User data. Gitignored. Never in version control. \\
\end{longtable}
}

\begin{center}\rule{0.5\linewidth}{0.5pt}\end{center}

\section{21. Shell Coding Standards}\label{shell-coding-standards}

Every script in the project must follow these rules. Violations are Gate
B failures.

\begin{Shaded}
\begin{Highlighting}[]
\CommentTok{\#!/usr/bin/env bash}
\BuiltInTok{set} \AttributeTok{{-}euo}\NormalTok{ pipefail}
\end{Highlighting}
\end{Shaded}

\begin{itemize}
\tightlist
\item
  \textbf{Bash, not sh.} \texttt{\#!/usr/bin/env\ bash} on every script.
  \texttt{\#!/bin/sh} is not acceptable.
\item
  \textbf{\texttt{set\ -euo\ pipefail}} is the second line. No
  exceptions.

  \begin{itemize}
  \tightlist
  \item
    \texttt{-e}: exit on error
  \item
    \texttt{-u}: exit on undefined variable
  \item
    \texttt{-o\ pipefail}: pipe failure propagates
  \end{itemize}
\item
  \textbf{Error propagation via return codes}, not exit traps or
  \texttt{exit\ 1} in lib functions.
\item
  \textbf{Logging via \texttt{lib/logging.sh}} functions. Never
  \texttt{echo} for user-facing messages in lib or services.
\item
  \textbf{Absolute paths always.} Use \texttt{\$WORKWARRIOR\_BASE/...},
  never relative paths.
\item
  \textbf{Quote all variable expansions:} \texttt{"\$var"} not
  \texttt{\$var}. \texttt{"\$\{array{[}@{]}\}"} not
  \texttt{\$\{array{[}@{]}\}}.
\item
  \textbf{Functions in snake\_case.} No camelCase. No kebab-case.
\item
  \textbf{All local variables declared with \texttt{local}.}
\item
  \textbf{No \texttt{cd} in lib functions.} Use full paths.
\item
  \textbf{Exit codes:} 0 = success, 1 = user error, 2 = system/internal
  error.
\end{itemize}

\subsection{Variable Naming}\label{variable-naming}

{\def\LTcaptype{none} % do not increment counter
\begin{longtable}[]{@{}ll@{}}
\toprule\noalign{}
Scope & Convention \\
\midrule\noalign{}
\endhead
\bottomrule\noalign{}
\endlastfoot
Global env vars & \texttt{UPPER\_SNAKE\_CASE} \\
Local function vars & \texttt{lower\_snake\_case} with \texttt{local} \\
Function names & \texttt{lower\_snake\_case} \\
\end{longtable}
}

\subsection{Function Template}\label{function-template}

\begin{Shaded}
\begin{Highlighting}[]
\FunctionTok{function\_name()} \KeywordTok{\{}
  \BuiltInTok{local} \VariableTok{arg1}\OperatorTok{=}\StringTok{"}\VariableTok{$\{1}\OperatorTok{:?}\NormalTok{function\_name requires arg1}\VariableTok{\}}\StringTok{"}
  \BuiltInTok{local} \VariableTok{arg2}\OperatorTok{=}\StringTok{"}\VariableTok{$\{2}\OperatorTok{:{-}}\NormalTok{default\_value}\VariableTok{\}}\StringTok{"}
  \BuiltInTok{local} \VariableTok{result}

  \CommentTok{\# implementation}

  \BuiltInTok{echo} \StringTok{"}\VariableTok{$result}\StringTok{"}
  \ControlFlowTok{return} \DecValTok{0}
\KeywordTok{\}}
\end{Highlighting}
\end{Shaded}

\begin{center}\rule{0.5\linewidth}{0.5pt}\end{center}

\section{22. Testing}\label{testing}

\subsection{BATS Tests}\label{bats-tests}

\begin{Shaded}
\begin{Highlighting}[]
\ExtensionTok{bats}\NormalTok{ tests/                          }\CommentTok{\# All BATS test suites}
\ExtensionTok{bats}\NormalTok{ tests/test{-}models{-}service.bats  }\CommentTok{\# Specific suite}
\ExtensionTok{bats}\NormalTok{ tests/test{-}profile{-}manager.bats}
\ExtensionTok{bats}\NormalTok{ tests/test{-}shell{-}integration.bats}
\ExtensionTok{bats}\NormalTok{ tests/test{-}github{-}sync.bats}
\end{Highlighting}
\end{Shaded}

BATS (Bash Automated Testing System) provides shell-native test
assertions. Each test file covers one service or lib component.

\subsection{Python Tests (pytest)}\label{python-tests-pytest}

\begin{Shaded}
\begin{Highlighting}[]
\ExtensionTok{python3} \AttributeTok{{-}m}\NormalTok{ pytest services/browser/  }\CommentTok{\# Browser server and heuristic engine}
\end{Highlighting}
\end{Shaded}

Tests cover: HTTP endpoint behavior, heuristic rule matching, compound
command splitting, AI fallback routing, SSE event emission.

\subsection{Test Policy}\label{test-policy}

\begin{itemize}
\tightlist
\item
  Tests run locally. No CI currently active.
\item
  Every implementation task (Gate C) requires test execution before
  merge.
\item
  Sync engine changes require integration tests against a test GitHub
  repository.
\item
  Shell function changes require BATS suite updates.
\end{itemize}

\begin{center}\rule{0.5\linewidth}{0.5pt}\end{center}

\section{23. Install Policy}\label{install-policy}

\subsection{Core Toolchain}\label{core-toolchain}

\texttt{ww\ deps\ install} is the canonical install path.

{\def\LTcaptype{none} % do not increment counter
\begin{longtable}[]{@{}lll@{}}
\toprule\noalign{}
Tool & Minimum Version & Package \\
\midrule\noalign{}
\endhead
\bottomrule\noalign{}
\endlastfoot
TaskWarrior & 3.0.0 & \texttt{task} \\
TimeWarrior & 1.4.0 & \texttt{timew} \\
JRNL & 4.0.0 & \texttt{jrnl} (via pipx) \\
Hledger & 1.30 & \texttt{hledger} \\
Bugwarrior & 1.22.0 & \texttt{bugwarrior} (via pipx) \\
GitHub CLI & 2.0.0 & \texttt{gh} \\
Python 3 & 3.9 & \texttt{python3} \\
pipx & any & \texttt{pipx} \\
\end{longtable}
}

\subsection{Platform Matrix}\label{platform-matrix}

{\def\LTcaptype{none} % do not increment counter
\begin{longtable}[]{@{}lll@{}}
\toprule\noalign{}
Platform & Manager & Auto-install \\
\midrule\noalign{}
\endhead
\bottomrule\noalign{}
\endlastfoot
macOS & brew & Yes \\
Debian/Ubuntu & apt & No (emits command) \\
Fedora/RHEL & dnf & No (emits command) \\
Arch/Manjaro & pacman & No (emits command) \\
Other & --- & Emits manual URLs \\
\end{longtable}
}

\subsection{Install Steps}\label{install-steps}

\begin{Shaded}
\begin{Highlighting}[]
\FunctionTok{git}\NormalTok{ clone }\OperatorTok{\textless{}}\NormalTok{repo{-}url}\OperatorTok{\textgreater{}}\NormalTok{ \textasciitilde{}/ww}
\BuiltInTok{cd}\NormalTok{ \textasciitilde{}/ww}
\ExtensionTok{./install.sh}
\BuiltInTok{source}\NormalTok{ \textasciitilde{}/.bashrc   }\CommentTok{\# or source \textasciitilde{}/.zshrc}
\ExtensionTok{ww}\NormalTok{ deps check}
\ExtensionTok{ww}\NormalTok{ profile create work}
\ExtensionTok{p{-}work}
\end{Highlighting}
\end{Shaded}

\subsection{Optional Extensions}\label{optional-extensions}

\begin{Shaded}
\begin{Highlighting}[]
\ExtensionTok{ww}\NormalTok{ tui install     }\CommentTok{\# taskwarrior{-}tui (full{-}screen TUI)}
\ExtensionTok{ww}\NormalTok{ mcp install     }\CommentTok{\# MCP server (AI agent access)}
\end{Highlighting}
\end{Shaded}

MCP install requires \texttt{uv} (Python package manager).
Auto-installed on macOS via brew.

\begin{center}\rule{0.5\linewidth}{0.5pt}\end{center}

\section{24. Dependency Table}\label{dependency-table}

{\def\LTcaptype{none} % do not increment counter
\begin{longtable}[]{@{}
  >{\raggedright\arraybackslash}p{(\linewidth - 6\tabcolsep) * \real{0.3529}}
  >{\raggedright\arraybackslash}p{(\linewidth - 6\tabcolsep) * \real{0.1765}}
  >{\raggedright\arraybackslash}p{(\linewidth - 6\tabcolsep) * \real{0.1765}}
  >{\raggedright\arraybackslash}p{(\linewidth - 6\tabcolsep) * \real{0.2941}}@{}}
\toprule\noalign{}
\begin{minipage}[b]{\linewidth}\raggedright
Dependency
\end{minipage} & \begin{minipage}[b]{\linewidth}\raggedright
Type
\end{minipage} & \begin{minipage}[b]{\linewidth}\raggedright
Role
\end{minipage} & \begin{minipage}[b]{\linewidth}\raggedright
Upstream
\end{minipage} \\
\midrule\noalign{}
\endhead
\bottomrule\noalign{}
\endlastfoot
TaskWarrior & Core tool & Task management, UDAs, hooks &
taskwarrior.org \\
TimeWarrior & Core tool & Time tracking & timewarrior.net \\
JRNL & Core tool & Journal entries & jrnl.sh \\
Hledger & Core tool & Double-entry accounting & hledger.org \\
Bugwarrior & Core tool & Issue pull from 20+ services &
github.com/GothenburgBitFactory/bugwarrior \\
GitHub CLI (gh) & Required & GitHub sync authentication and API &
cli.github.com \\
Python 3 (≥3.9) & Required & Browser UI server & python.org \\
pipx & Required & JRNL and Bugwarrior isolation & pypa.github.io/pipx \\
taskgun & Weapon & Bulk task series (Gun weapon) &
github.com/hamzamohdzubair/taskgun \\
taskcheck & Weapon & Auto-scheduler (Schedule weapon) & --- \\
taskwarrior-tui & Optional & Full-screen TUI &
github.com/kdheepak/taskwarrior-tui \\
taskwarrior-mcp & Optional & MCP server for AI agents &
github.com/hnsstrk/taskwarrior-mcp \\
ollama & Optional AI & Local LLM provider & ollama.ai \\
TWDensity & Optional extension & Due-date density scoring & --- \\
\end{longtable}
}

\begin{center}\rule{0.5\linewidth}{0.5pt}\end{center}

\emph{Workwarrior Technical Standard v0.1.0 --- Babb ---
workwarrior@babb.tel}

\end{document}
